\originalchapter{8.05 作业3：自旋旋转与换基}
来源为Fall 2013 Assignment 3, 2013年9月20日发布, 9月27日15时截止. 本篇处理第1--4、5(b)(c)、6题, 解答均为\textbf{AI独立解答}. 第5(a)明确打印Axler题面, 该具体外书对象暂列引用边界, 不影响本作业其余内容.
\originalsection{第1题：方向与连续筛选}
\begin{exercise}
(a)未归一态$(1+\ii)\ket+-(1+\ii\sqrt3)\ket-$指哪个方向? (b)先制备$z+$, 再沿$xz$面极角$\theta$测量并只留$n+$, 最后测$z$. 相对进入第二磁铁总数, 求最后$z+$、$z-$及未到第三磁铁比例, 核$\theta=0,\pi/2,\pi$.
\end{exercise}
\begin{solution}
模平方和为6. 对系数$a,b$, Bloch向量为$(2\operatorname{Re}a^*b,2\operatorname{Im}a^*b,|a|^2-|b|^2)/6$. 实际原PDF括号内是$-(1+\ii\sqrt3)$, 得$\boldsymbol n=(-(1+\sqrt3)/3,(1-\sqrt3)/3,-1/3)$, 单位范数可直接核. 因而$\theta=\arccos(-1/3)$, $\phi=13\pi/12$模$2\pi$.

第二磁铁保留概率$c^2=\cos^2(\theta/2)$, 保留后的态为$n+$, 第三次条件$z+$概率$c^2$、$z-$概率$s^2=\sin^2(\theta/2)$. 所以相对原束比例为$c^4,c^2s^2,s^2$, 三者和1. $\theta=0$得$(1,0,0)$, $\pi/2$得$(1/4,1/4,1/2)$, $\pi$得$(0,0,1)$, 后者是第二次全被挡住, 不能将条件概率误作总束概率.
\end{solution}
\originalsection{第2题：半角重叠}
\begin{exercise}
两自旋$n+,n'+$方向的夹角为$\gamma$, 直接证$|\braket{n';+}{n;+}|^2=\cos^2(\gamma/2)$.
\end{exercise}
\begin{solution}
写$\ket{n;+}=(\cos\frac\theta2,\e^{\ii\phi}\sin\frac\theta2)^{\mathsf T}$, 内积为$c'c+s's\e^{\ii(\phi-\phi')}$. 模平方用半角恒等式整理为$[1+\cos\theta\cos\theta'+\sin\theta\sin\theta'\cos(\phi-\phi')]/2=(1+\boldsymbol n\cdot\boldsymbol n')/2$. 因$\boldsymbol n\cdot\boldsymbol n'=\cos\gamma$, 即为半角平方. 两方向相反而态正交, 同方向虽可能整体相位不同却重叠模1.
\end{solution}
\originalsection{第3题：旋转算符}
\begin{exercise}
$R_n(\alpha)=\e^{-\ii\alpha\boldsymbol n\cdot\boldsymbol\sigma/2}$. 从级数求显式式并证幺正; 求$R_yS_zR_y^\dagger$及$R_y\ket+$, 它是什么分量的正本征态?
\end{exercise}
\begin{solution}
Pauli乘积给$(\boldsymbol n\cdot\sigma)^2=I$, 分组奇偶幂得$R_n=I\cos(\alpha/2)-\ii\boldsymbol n\cdot\sigma\sin(\alpha/2)$. 厄米性与平方$I$使$R_n^\dagger R_n=I$. $R_y=\begin{pmatrix}c&-s\\s&c\end{pmatrix}$, 直接乘得$R_yS_zR_y^\dagger=S_z\cos\alpha+S_x\sin\alpha$. $R_y\ket+=c\ket++s\ket-$, 正是方向$(\sin\alpha,0,\cos\alpha)$的$+\hbar/2$本征态. 原题用$\hbar=1$记号, 此处恢复单位. $R(2\pi)=-I$使态整体变号, 同一个自旋方向不变.
\end{solution}
\originalsection{第4题：Schwarz与三角不等式}
\begin{exercise}
以$\norm{a-\lambda b}^2\ge0$分别在实、复空间证Schwarz并给等号; 再证三角不等式及等号条件.
\end{exercise}
\begin{solution}
$b\ne0$时复空间展开为$\norm a^2-\lambda\braket a b-\lambda^*\braket b a+|\lambda|^2\norm b^2$. 极小$\lambda=\braket b a/\norm b^2$给$\norm a^2-|\braket b a|^2/\norm b^2\ge0$. 实数情形只将共轭删去. 等号恰$a=\lambda b$, 也包括任一零向量. 再展开$\norm{a+b}^2=\norm a^2+\norm b^2+2\operatorname{Re}\braket a b\le(\norm a+\norm b)^2$. 非零两矢时同时要求共线与内积为正实, 即$a=cb,c>0$. 若允许零矢, 也有等号; 原题“当且仅当正数倍”漏了零矢的边界, 此处明确补全.
\end{solution}
\originalsection{第5题：无限维与单射}
\begin{exercise}
只处理(b)(c): 无限序列空间$F^\infty$为何无限维? 线性映射$T$单射当且仅当$\ker T=\{0\}$.
\end{exercise}
\begin{solution}
对每个正整数$N$, 前$N$个标准序列$e_j$在第$j$位置为1、其他0, 线性独立; 若存在有限维$d$, 取$N=d+1$矛盾. 这是维数无限, 不要求每个无限序列都由有限线性组合的$e_j$组成. 单射使$Tv=T0$只能$v=0$. 反向若$Tu=Tv$, 则$T(u-v)=0$, 零核推出$u=v$. (a)外书明确题面尚单核, 不把它记为已解.
\end{solution}
\originalsection{第6题：换基不变量}
\begin{exercise}
$u_j=Av_j$, $B=A^{-1}$, 以两基的展开证明相应换基矩阵互逆. 对线性算符$T$求两基矩阵关系, 证明迹和行列式不依基.
\end{exercise}
\begin{solution}
设$C$的第$j$列为$u_j$在$v$基坐标, 则$[w]_v=C[w]_u$. 反向展开$[w]_u=D[w]_v$, 对任意$w$有$DC=CD=I$, 故$D=C^{-1}$. 将$Tw$两边表示为$T_vC[w]_u=CT_u[w]_u$, 得$T_u=C^{-1}T_vC$. 这里$C$亦为$A$在$v$基的矩阵, $D$为$B$在$u$基的矩阵. 迹循环性由$\Tr(XY)=\sum_{ij}X_{ij}Y_{ji}=\Tr(YX)$直接证, 所以$\Tr T_u=\Tr T_v$; 行列式乘法则给$\det T_u=\det C^{-1}\det T_v\det C=\det T_v$. 基无须正交; 非正交换基要用逆矩阵而非随意用伴随.
\end{solution}
